\daytitle{13}{机械振动}

\begin{learningbox}{本节导读}
机械振动是自然界最普遍的周期性运动形式。本节学习简谐运动的基本特征——弹簧振子和单摆，理解回复力、振幅、周期、频率等概念，以及简谐运动的 $x$-$t$ 图像。
\end{learningbox}

\textbf{本节读题路线：} 先找平衡位置，再从平衡位置量位移 $x$；接着判断回复力是否与位移成正比并反向。图像题先读振幅、周期和初相位；周期题先判断是弹簧振子、单摆，还是需要先做等效处理。

\section*{一、学习目标}

\begin{enumerate}
  \item 能用回复力 $F = -Kx$ 判断一个运动是否可近似看作简谐运动，并区分一般回复系数 $K$ 与弹簧劲度系数 $k$。
  \item 能说明弹簧振子周期 $T = 2\pi\sqrt{\dfrac{m}{k}}$ 与惯性大小、回复强弱的关系。
  \item 能说明理想单摆周期 $T = 2\pi\sqrt{\dfrac{l}{g}}$ 的小角度适用条件。
  \item 能读懂简谐运动的 $x$-$t$ 图像，知道 $A$、$T$、$\omega$ 的物理意义。
  \item 能结合图像说明简谐运动中的动能与势能转化。
  \item 能把多个弹簧先等效为一个弹簧，再代入弹簧振子模型。
\end{enumerate}
\section*{二、初高中衔接}

初中物理通过“弹簧振子”和“单摆”定性接触过振动——知道摆来回摆动一次是一周期。

高中把振动定量化：定义了振幅 $A$、周期 $T$、频率 $f$，并给出了弹簧振子和单摆的周期公式——这是中学生第一次可以精确计算振动系统的周期。

\textbf{科学发展视角：} 振动既可以被利用，也需要被控制。石英钟靠稳定振动计时，建筑和桥梁要避开共振，手机和汽车中的振动传感器能感知运动状态。学习简谐运动时，要把“回复力指向平衡位置”看成很多复杂振动的最基本模型。

\section*{三、核心概念}

\begin{conceptbox}{简谐运动}

\textbf{定义：} 物体所受回复力（或沿振动方向的合力）与偏离平衡位置的位移大小成正比，方向总是指向平衡位置，这样的运动叫简谐运动。

\textbf{回复力：}
\[
\boxed{F = -Kx}
\]

\begin{itemize}
  \item $K$：回复力大小与位移大小的比例系数，单位为 $\mathrm{N/m}$；它由具体振动系统决定。
  \item $x$：偏离平衡位置的位移。
  \item 负号表示回复力方向始终与位移方向相反（指向平衡位置）。
\end{itemize}

\textbf{符号说明：} 对水平弹簧振子，回复系数 $K$ 就是弹簧的劲度系数 $k$；对其他简谐运动模型，$K$ 不一定表示某根弹簧的劲度系数。

\textbf{描述振动的物理量：}
\begin{itemize}
  \item \textbf{振幅 $A$：} 最大位移的绝对值，反映振动强弱。
  \item \textbf{周期 $T$：} 完成一次全振动所需的时间。
  \item \textbf{频率 $f$：} 单位时间内完成全振动的次数，$f = \dfrac{1}{T}$。
\end{itemize}

\end{conceptbox}

\begin{conceptbox}{弹簧振子}

水平弹簧振子是最基本的简谐运动模型。

\begin{center}
% TikZ 图：弹簧振子
% Day 13：机械振动
\begin{tikzpicture}[scale=0.9,>=Latex]
  % 墙壁
  \draw[thick] (-2,-0.5) -- (-2,1.5);
  % 弹簧（平衡位置）
  \draw[decorate,decoration={coil,aspect=0.4,segment length=3pt,amplitude=3pt}] (-1.5,0.5) -- (0.5,0.5);
  \draw[thick] (0.5,0) rectangle (1.8,1);
  \node at (1.15,0.5) {$m$};
  \node[above] at (-0.5,0.5) {平衡位置};
  \node[below] at (1.15,-0.2) {$x=0$};
  % 最大位移位置（虚线）
  \draw[dashed] (3.5,0) -- (3.5,1);
  \node[above] at (3.5,1) {$x=A$};
  % 回复力
  \draw[->,thick,red] (1.8,0.5) -- (3.3,0.5) node[above] {$F$};
  % 位移
  \draw[<->] (1.15,-0.3) -- (3.5,-0.3) node[midway,below] {$x$};
\end{tikzpicture}


\figcaption{弹簧振子：$F = -kx$，$x$ 是偏离平衡位置的位移}\label{fig:spring-oscillator}
\end{center}

\textbf{周期公式：}
\[
\boxed{T = 2\pi\sqrt{\frac{m}{k}}}
\]

\textbf{公式来源：}
弹簧振子偏离平衡位置 $x$ 时，合力为
\[
F=-kx.
\]
把回复力代入牛顿第二定律 $F=ma$，得到
\[
a=-\frac{k}{m}x.
\]
把加速度写成位移对时间的二阶变化，运动方程为
\[
\frac{\mathrm d^2x}{\mathrm dt^2}+\frac{k}{m}x=0.
\]
不要求在这里求解这个微分方程。把 $x=A\cos(\omega t+\varphi_0)$ 代入上式可以验证：它能成立的条件是
\[
\omega^2=\frac{k}{m},\qquad T=\frac{2\pi}{\omega}=2\pi\sqrt{\frac{m}{k}}.
\]
负号表示加速度总指向平衡位置；$\dfrac{k}{m}$ 决定振动的快慢。因此弹簧越硬（$k$ 越大）振得越快，物体质量越大（$m$ 越大）振得越慢。

\textbf{学习边界：} 这里不要求会求解微分方程。对一维、质量不变、阻尼可忽略的理想模型，回复力满足 $F=-Kx$ 时，系统做简谐运动；周期由惯性大小和回复强弱共同决定。真实振动若阻尼明显、回复力不再与位移成正比，就不能直接套用简谐运动结论。

\begin{itemize}
  \item $m$：振子质量。
  \item $k$：弹簧劲度系数。
  \item 在线性弹簧模型和无阻尼近似下，周期与振幅无关（简谐运动的等时性）。
\end{itemize}

\end{conceptbox}

\begin{conceptbox}{单摆}

在摆角很小时，单摆的运动可\textbf{近似}看作简谐运动。中学测量实验通常把摆角控制在 $5^\circ$ 以内；可接受的最大摆角取决于实验精度要求，$5^\circ$ 不是绝对分界。

\begin{center}
% TikZ 图：单摆
% Day 13：机械振动
\begin{tikzpicture}[scale=0.9,>=Latex]
  % 天花板
  \draw[thick] (-1.5,0) -- (1.5,0);
  \filldraw[black] (0,0) circle (2pt);
  % 摆线（最大振幅位置）
  \draw[thick] (0,0) -- (-1.2,-2.5);
  \filldraw[red] (-1.2,-2.5) circle (3pt);
  \node[left] at (-1.2,-2.5) {$m$};
  % 平衡位置（虚线）
  \draw[dashed] (0,0) -- (0,-2.5);
  \filldraw[black] (0,-2.5) circle (2pt);
  \node[right] at (0,-2.5) {平衡};
  % 摆角
  \draw (0,-0.4) arc (-90:-120:0.4);
  \node at (-0.3,-0.7) {$\theta$};
  % 回复力示意
  \draw[->,thick,red] (-1.2,-2.5) -- ++(0.8,0) node[right] {$F = -mg\sin\theta$};
  % 摆长
  \draw[<->,thick] (-1.5,0) -- (-1.5,-2.5) node[midway,left] {$l$};
\end{tikzpicture}


\figcaption{单摆：以沿圆弧方向、偏离平衡位置的弧长为 $x$，小角度时 $F_\tau\approx-\dfrac{mg}{l}x$}\label{fig:pendulum}
\end{center}

\textbf{周期公式：}
\[
\boxed{T = 2\pi\sqrt{\frac{l}{g}}}
\]

\textbf{公式来源（小角度近似）：}
摆球沿圆弧方向的回复力来自重力的切向分力：
\[
F_{\tau}=-mg\sin\theta.
\]
若用弧长 $x$ 描述摆球偏离平衡位置的程度，角度用弧度表示时有准确关系 $x=l\theta$。当摆角很小时，近似关系 $\sin\theta\approx\theta$ 成立，所以
\[
F_{\tau}\approx-mg\theta=-\frac{mg}{l}x.
\]
它与弹簧振子的 $F=-kx$ 在数学形式上相同，只是等效“劲度系数”为
\[
k_{\text{等效}}=\frac{mg}{l}.
\]
代入 $T=2\pi\sqrt{\dfrac{m}{k}}$，得
\[
T=2\pi\sqrt{\frac{m}{mg/l}}=2\pi\sqrt{\frac{l}{g}}.
\]
这也解释了单摆周期为什么与摆球质量无关。

\begin{itemize}
  \item $l$：摆长，指悬点到摆球重心的距离。摆球可视为均匀小球时，其重心在球心，因此摆长应量到球心，不能只把绳长作为摆长。
  \item $g$：当地重力加速度。
  \item 在理想单摆、小角度近似下，周期与摆球质量无关，并近似与振幅无关——这是惠更斯的重要发现。
  \item 利用单摆可以测量当地的重力加速度。
\end{itemize}

\end{conceptbox}

\begin{conceptbox}{简谐运动的 $x$-$t$ 图像}

\begin{center}
% TikZ 图：简谐运动的 x-t 图像（正弦曲线）
% Day 13：机械振动
\begin{tikzpicture}[scale=1,>=Latex]
  \draw[->] (-0.3,0) -- (6.5,0) node[right] {$t$};
  \draw[->] (0,-1.8) -- (0,1.8) node[above] {$x$};
  \node[below left] at (0,0) {$O$};
  % 正弦曲线
  \draw[very thick,red,smooth] (0,0) sin (1.5,1.5) cos (3,0) sin (4.5,-1.5) cos (6,0);
  % 振幅标注
  \draw[dashed] (0,1.5) -- (1.5,1.5);
  \draw[<->] (0.35,0) -- (0.35,1.5) node[midway,right] {$A$};
  % 周期标注：从同相位点到下一个同相位点
  \draw[<->] (0,-1.0) -- (6,-1.0);
  \node[below] at (3,-1.0) {$T$};
  % 公式
  \node at (3,-2.2) {$x = A\cos(\omega t + \varphi_0)$};
\end{tikzpicture}


\figcaption{简谐运动的位移-时间图像为正弦（或余弦）曲线}\label{fig:sh-motion-graph}
\end{center}

图中选择 $t=0$ 时振子位于正向最大位移处，因此初相位为 $\varphi_0=0$，表达式为
\[
x=A\cos\omega t.
\]
若计时起点不同，更一般地写成
\[
x=A\cos(\omega t+\varphi_0),
\]
其中 $\omega=\dfrac{2\pi}{T}$ 为角频率，$\varphi_0$ 描述开始计时时振子处于振动过程中的哪个状态。

\textbf{由位移得到速度和加速度：} 位移对时间的变化率是速度，速度对时间的变化率是加速度，因此
\[
v=\frac{\mathrm dx}{\mathrm dt}
  =-A\omega\sin(\omega t+\varphi_0),\qquad
a=\frac{\mathrm dv}{\mathrm dt}
  =-\omega^2x.
\]
由 $|\sin(\omega t+\varphi_0)|\leqslant 1$ 可得
\[
\boxed{v_{\max}=\omega A}.
\]
初学阶段不要求独立完成求导，但要会使用图像斜率判断速度，并理解：在平衡位置速度大小最大，在两端速度为零。

\textbf{读图要点：} 先从图上读出振幅 $A$、周期 $T$ 和 $t=0$ 时的位置；再看曲线上该点切线的斜率，斜率就是瞬时速度 $v=\dfrac{\mathrm dx}{\mathrm dt}$。

\end{conceptbox}

\section*{四、公式与适用条件}

\textbf{模型使用前检查：}
\begin{enumerate}
  \item 先找平衡位置，位移 $x$ 必须从平衡位置量起。
  \item 判断回复力是否与位移成正比且方向相反，这是简谐运动的核心。
  \item 单摆周期公式只适用于小角度近似，摆长应量到摆球球心。
  \item 区分质点振动速度和波速，振动章节中的 $v$ 要看语境。
\end{enumerate}

\begin{center}
\begin{tabular}{lll}
\toprule
公式 & 物理意义 & 适用条件 \\
\midrule
$F = -Kx$ & 简谐运动的回复力 & 位移从平衡位置计，线性回复力 \\
$T = 2\pi\sqrt{\dfrac{m}{k}}$ & 弹簧振子周期 & 轻弹簧、线性范围、忽略阻尼 \\
$T = 2\pi\sqrt{\dfrac{l}{g}}$ & 单摆周期 & 理想单摆，小角度近似 \\
$x = A\cos(\omega t + \varphi_0)$ & 振动方程 & 简谐运动 \\
$f = \dfrac{1}{T}$ & 频率与周期 & 任何周期性运动 \\
\bottomrule
\end{tabular}
\end{center}

\begin{learningbox}{简谐运动的状态分析清单}
判断简谐运动时，先把“位置、速度、加速度、能量”对应起来，比单独记公式更可靠。
\begin{center}
\small
\begin{tabular}{p{0.17\textwidth}p{0.18\textwidth}p{0.22\textwidth}p{0.23\textwidth}}
\toprule
位置 & 速度 & 加速度/回复力 & 能量特点 \\
\midrule
$x=0$ 平衡位置 & 速度大小最大 & $a=0$，回复力为零 & 动能最大，势能最小 \\
$x=\pm A$ 端点 & 瞬时速度为零 & 加速度大小最大，方向指向平衡位置 & 势能最大，动能为零 \\
从端点向平衡位置运动 & 速度增大 & 加速度方向指向平衡位置 & 势能转化为动能 \\
从平衡位置向端点运动 & 速度减小 & 加速度方向仍指向平衡位置 & 动能转化为势能 \\
\bottomrule
\end{tabular}
\end{center}
\textbf{读图提醒：} $x$-$t$ 图像上某点切线的斜率就是该时刻的瞬时速度；曲线离平衡位置越远，回复力和加速度的大小越大。
\end{learningbox}

\section*{五、典型模型}

\subsection*{模型 1：水平弹簧振子}

平衡位置处速度最大、加速度为零；最大位移处速度为零、加速度大小最大。
理想无阻尼弹簧振子机械能守恒：
\[
\dfrac12 kA^2 = \dfrac12 mv^2 + \dfrac12 kx^2.
\]

\subsection*{模型 2：单摆测 $g$}

由 $T = 2\pi\sqrt{\dfrac{l}{g}}$ 得 $g = \dfrac{4\pi^2 l}{T^2}$。实验中改变摆长测周期，作图求斜率可得 $g$。

更规范的数据处理方法是把公式平方：
\[
T^2=\frac{4\pi^2}{g}l.
\]
以 $l$ 为横轴、$T^2$ 为纵轴作图，理想关系是一条过原点的直线；实验数据有随机误差时，各点会分布在最佳拟合直线附近。若拟合直线斜率为 $s$，则
\[
s=\frac{4\pi^2}{g},\qquad [s]=\mathrm{s^2/m},\qquad
g=\frac{4\pi^2}{s}.
\]
若最佳拟合直线具有明显非零截距，应检查摆长零点、计时起止位置等系统误差，不能只凭任意两点连线求斜率。

\begin{center}
% TikZ 图：单摆实验数据处理
% Day 13：机械振动
\begin{tikzpicture}[scale=0.82,>=Latex]
  \draw[->,thick] (0,0) -- (4.9,0) node[right] {$l/\mathrm{m}$};
  \draw[->,thick] (0,0) -- (0,3.45) node[above] {$T^2/\mathrm{s^2}$};

  % 理想关系通过原点，实验点在最佳拟合直线附近波动。
  \draw[blue!65,thick] (0,0) -- (4.45,2.85);
  \node[blue!70,anchor=west] at (0.55,3.15)
    {$T^2=\dfrac{4\pi^2}{g}l$};

  \foreach \x/\y in {0.80/0.50,1.45/0.96,2.15/1.35,2.85/1.86,3.55/2.24,4.15/2.69}{
    \filldraw[red] (\x,\y) circle (2pt);
  }

  \draw[dashed,gray] (2.25,0) -- (2.25,1.44) -- (0,1.44);
  \draw[dashed,gray] (3.75,0) -- (3.75,2.40) -- (0,2.40);
  \draw[<->] (2.25,0.28) -- (3.75,0.28);
  \node[below] at (3.0,0.17) {$\Delta l$};
  \draw[<->] (4.05,1.44) -- (4.05,2.40);
  \node[right] at (4.10,1.92) {$\Delta T^2$};

  \draw[rounded corners,gray] (5.10,0.35) rectangle (7.70,2.90);
  \node[align=left,font=\small,anchor=west] at (5.32,2.35)
    {斜率关系：};
  \node[align=left,font=\small,anchor=west] at (5.32,1.90)
    {$s=\Delta T^2/\Delta l$};
  \node[align=left,font=\small,anchor=west] at (5.32,1.42)
    {$s=4\pi^2/g$};
  \node[align=left,font=\small,anchor=west] at (5.32,0.92)
    {$[s]=\mathrm{s^2/m}$};
  \node[align=left,font=\small,anchor=west] at (5.32,0.52)
    {$g=4\pi^2/s$};
\end{tikzpicture}


\figcaption{单摆测 $g$ 的数据处理：作 $T^2-l$ 图像，用斜率求 $g$}\label{fig:pendulum-data-plot}
\end{center}

\subsection*{模型 3：简谐运动的能量}

对满足 $F=-Kx$ 的理想线性简谐模型，总机械能可写为
\[
E=\frac12 KA^2.
\]
对水平弹簧振子，$K=k$；对单摆小角度振动，$K=mg/l$。动能和势能不断相互转化，总和守恒。

\begin{center}
% TikZ 图：简谐运动能量随位置变化
% Day 13：机械振动
\begin{tikzpicture}[scale=0.92,>=Latex]
  \draw[->] (-3.15,0) -- (3.35,0) node[right] {$x$};
  \draw[->] (0,0) -- (0,2.75) node[above] {能量};

  \draw[dashed] (-2.65,2.20) -- (2.65,2.20);
  \node[right] at (2.65,2.20) {$E=\dfrac12 KA^2$};

  \draw[very thick,red,smooth,domain=-2.5:2.5,samples=80]
    plot (\x,{0.352*\x*\x});
  \draw[very thick,blue,smooth,domain=-2.5:2.5,samples=80]
    plot (\x,{2.20-0.352*\x*\x});

  \node[red] at (2.15,1.88) {$E_p$};
  \node[blue] at (1.05,1.82) {$E_k$};

  \draw[dashed] (-2.5,0) -- (-2.5,2.20);
  \draw[dashed] (2.5,0) -- (2.5,2.20);
  \node[below] at (-2.5,0) {$-A$};
  \node[below left] at (0,0) {$O$};
  \node[below] at (2.5,0) {$A$};

  \node[align=center] at (0,-0.82)
    {平衡位置动能最大；两端势能最大};
\end{tikzpicture}


\figcaption{简谐运动能量转化：动能和势能此消彼长，总能量守恒}\label{fig:shm-energy}
\end{center}

\section*{六、例题精讲}

\begin{examplebox}{弹簧振子的周期}

\textbf{题意分析：} 一个水平弹簧振子，$m = 0.5\ \mathrm{kg}$，$k = 50\ \mathrm{N/m}$。求：
\begin{enumerate}
  \item 振子的周期和频率。
  \item 若振幅 $A = 0.1\ \mathrm{m}$，求系统的总机械能。
\end{enumerate}

\textbf{物理模型：} 弹簧振子简谐运动。

\textbf{解题过程：}

\begin{enumerate}
  \item 周期：
  \[
  T = 2\pi\sqrt{\frac{m}{k}} = 2\pi\sqrt{\frac{0.5}{50}} = 2\pi\sqrt{0.01} = 2\pi \times 0.1 \approx 0.628\ \mathrm{s}
  \]
  频率：
  \[
  f = \frac{1}{T} \approx 1.59\ \mathrm{Hz}
  \]

  \item 总机械能等于最大弹性势能：
  \[
  E = \frac12 k A^2 = \frac12 \times 50 \times 0.1^2 = 0.25\ \mathrm{J}
  \]
\end{enumerate}

\textbf{方法总结：} 弹簧振子题先做条件检查：轻弹簧、线性范围、阻尼可忽略。满足这些条件时，周期只由 $m$ 和 $k$ 决定；若题目问能量，再看是否给出振幅或最大速度，分别用 $\frac12 kA^2$ 或 $\frac12 mv_{\max}^2$。

\end{examplebox}

\begin{examplebox}{单摆测重力加速度}

\textbf{题意分析：} 某同学用单摆测重力加速度，测得摆长 $l = 0.992\ \mathrm{m}$，50 次全振动的时间 $t = 99.5\ \mathrm{s}$。求重力加速度 $g$。

\textbf{物理模型：} 单摆实验数据处理。

\textbf{解题过程：}

周期 $T = \dfrac{t}{50} = \dfrac{99.5}{50} = 1.99\ \mathrm{s}$。

由 $T = 2\pi\sqrt{\dfrac{l}{g}}$：
\[
g = \frac{4\pi^2 l}{T^2} = \frac{4 \times 3.14^2 \times 0.992}{1.99^2} \approx \frac{4 \times 9.86 \times 0.992}{3.96} \approx 9.88\ \mathrm{m/s^2}
\]

\textbf{方法总结：} 实验测 $g$ 时按三步检查：先数清全振动次数并求平均周期 $T=t/n$；再从悬点量到球心得到摆长 $l$；最后代入 $g=4\pi^2l/T^2$。若把摆长量到球面或只测一次周期，误差会明显放大。

\end{examplebox}

\begin{examplebox}{由振动状态判断速度、加速度和能量}

\textbf{题意分析：} 一个水平弹簧振子从正向最大位移处由静止释放，满足
\[
x=A\cos\frac{2\pi}{T}t.
\]
求 $t=\dfrac{T}{3}$ 时的位移、速度方向、加速度方向，以及动能和总机械能之比。

\textbf{物理模型：} 简谐运动的图像、状态与能量对应。

\textbf{解题过程：} 代入 $t=\dfrac{T}{3}$，
\[
x=A\cos\frac{2\pi}{3}=-\frac{A}{2}.
\]
此时振子位于平衡位置左侧。由
\[
v=-A\omega\sin\omega t,\qquad a=-\omega^2x
\]
可知 $v<0$，振子仍向负方向运动；$a>0$，加速度指向平衡位置。弹性势能为
\[
E_p=\frac12kx^2
   =\frac14\left(\frac12kA^2\right)
   =\frac14E,
\]
所以
\[
\boxed{\frac{E_k}{E}=1-\frac{E_p}{E}=\frac34}.
\]

\textbf{方法总结：} 先由 $x$ 确定位置，再由图像斜率或速度表达式判断运动方向；加速度始终满足 $a=-\omega^2x$。能量问题可直接使用 $E_p/E=x^2/A^2$。

\end{examplebox}

\section*{七、变式训练}

\begin{enumerate}
  \item 弹簧振子 $m = 0.2\ \mathrm{kg}$，$k = 80\ \mathrm{N/m}$，振幅 $A = 0.05\ \mathrm{m}$。求周期、频率和最大速度。

  \item 单摆摆长 $l = 2\ \mathrm{m}$，$g = 9.8\ \mathrm{m/s^2}$。求周期。若想让周期变为 $3\ \mathrm{s}$，应如何调整摆长？
  \begin{center}
  \begin{tikzpicture}[scale=0.9,>=Stealth,line cap=round,line join=round]
  \begin{scope}
    \path[use as bounding box] (-0.3,-1.05) rectangle (6.5,1.9);
    \draw[thick] (-0.2,0) -- (3.0,0);
    \draw[decorate,decoration={zigzag,segment length=5pt,amplitude=2pt},thick]
      (0.2,0.55) -- (1.35,0.55);
    \draw[fill=gray!12,thick] (1.35,0.22) rectangle (2.0,0.88);
    \draw[dashed] (1.675,0.05) -- (1.675,1.08);
    \draw[<->,blue,thick] (1.675,1.18) -- (2.55,1.18) node[midway,above] {$A$};
    \node[font=\small] at (1.5,-0.36) {弹簧振子};
    \node[font=\small] at (1.5,-0.76) {$T=2\pi\sqrt{m/k}$};
    \node[font=\small] at (0.55,0.92) {$k$};
    \node[font=\small] at (1.68,0.55) {$m$};
  \end{scope}
  \begin{scope}[xshift=4.95cm]
    \draw[thick] (-0.7,1.55) -- (0.7,1.55);
    \coordinate (O) at (0,1.55);
    \coordinate (M) at (0.55,0.15);
    \draw[thick] (O) -- (M);
    \draw[fill=gray!20,thick] (M) circle (0.13);
    \draw[dashed] (O) -- (0,-0.1);
    \draw[->] (0,-0.05) arc[start angle=-90,end angle=-69,radius=1.6];
    \node[font=\small] at (0.23,0.95) {$\theta$};
    \draw[<->,blue,thick] (-0.11,1.505) -- (0.44,0.105);
    \node[blue,font=\small] at (0.04,0.82) {$l$};
    \node[font=\small] at (0,-0.36) {单摆};
    \node[font=\small] at (0,-0.76) {$T=2\pi\sqrt{l/g}$};
  \end{scope}
\end{tikzpicture}

  \end{center}

  \item 某同学用单摆测重力加速度，摆长 $l = 1.0\ \mathrm{m}$，测得 40 次全振动的时间 $t = 80.0\ \mathrm{s}$。求 $g$。

  \item 两根轻弹簧劲度系数分别为 $k_1$、$k_2$。若并联后连接同一物体做水平振动，等效劲度系数是多少？若串联后连接同一物体，等效劲度系数满足什么关系？
\end{enumerate}

\begin{learningbox}{变式训练参考结果}
\begin{enumerate}[itemsep=0.45em,topsep=0.3em,parsep=0pt]
  \item $\omega=\sqrt{k/m}=20\ \mathrm{rad/s}$，
  $T=2\pi/\omega\approx0.314\ \mathrm s$，
  $f=1/T\approx3.18\ \mathrm{Hz}$，$v_{\max}=\omega A=1.0\ \mathrm{m/s}$。
  \item 原周期 $T=2\pi\sqrt{l/g}\approx2.84\ \mathrm s$。要使周期为
  $3.0\ \mathrm s$，应将摆长增大到
  $l'=g(3/2\pi)^2\approx2.23\ \mathrm m$。
  \item 单摆周期为 $T=t/40=2.00\ \mathrm s$，故
  $g=4\pi^2l/T^2=\pi^2\ \mathrm{m/s^2}\approx9.87\ \mathrm{m/s^2}$。
  \item 并联时 $k_{\text{并}}=k_1+k_2$；串联时
  $\displaystyle 1/k_{\text{串}}=1/k_1+1/k_2$。
\end{enumerate}
\end{learningbox}

\section*{八、竞赛入门补充}

\begin{competitionbox}
\textbf{简谐运动的证明}

判断一个运动是否是简谐运动的标准：以平衡位置为位移原点，看振动方向上的合力是否满足 $F=-Kx$。

\textbf{例题：} 竖直方向的弹簧振子。弹簧原长 $l_0$，挂上质量为 $m$ 的物体后伸长 $\Delta l$ 静止（平衡位置）。证明物体在竖直方向做简谐运动。

\begin{center}
\begin{tikzpicture}[scale=0.92, >=Stealth, every node/.style={font=\small}]
  % Left state: equilibrium
  \draw[thick] (-3.2,2.1) -- (-0.8,2.1);
  \foreach \x in {-3.1,-2.85,...,-0.95} {
    \draw[gray] (\x,2.1) -- ++(0.15,0.18);
  }
  \draw[decorate, decoration={coil, amplitude=3pt, segment length=5pt}, thick]
    (-2,2.1) -- (-2,0.65);
  \draw[thick] (-2.4,0.15) rectangle (-1.6,0.65);
  \node at (-2,0.40) {$m$};
  \draw[dashed] (-3.35,0.65) -- (-0.65,0.65);
  \node[anchor=east] at (-3.35,0.65) {平衡位置};
  \draw[->, blue, thick] (-1.45,0.40) -- (-1.45,-0.40) node[right] {$mg$};
  \draw[->, red, thick] (-2.55,0.65) -- (-2.55,1.45) node[left] {$F_{\text{弹}}$};
  \node[align=center] at (-2,-0.95) {平衡时\\$k\Delta l=mg$};

  % Right state: pulled down by x
  \draw[thick] (0.8,2.1) -- (3.2,2.1);
  \foreach \x in {0.9,1.15,...,3.05} {
    \draw[gray] (\x,2.1) -- ++(0.15,0.18);
  }
  \draw[decorate, decoration={coil, amplitude=3pt, segment length=5pt}, thick]
    (2,2.1) -- (2,-0.15);
  \draw[thick] (1.6,-0.65) rectangle (2.4,-0.15);
  \node at (2,-0.40) {$m$};
  \draw[dashed] (0.65,0.65) -- (3.35,0.65);
  \node[anchor=west] at (3.35,0.65) {平衡位置};
  \draw[<->, thick] (3.0,0.65) -- (3.0,-0.15)
    node[midway, right] {$x$};
  \draw[->, blue, thick] (2.55,-0.40) -- (2.55,-1.20) node[right] {$mg$};
  \draw[->, red, thick] (1.45,-0.15) -- (1.45,0.95)
    node[left] {$k(\Delta l+x)$};
  \node[align=center] at (2,-1.75) {下拉 $x$ 后\\$F=mg-k(\Delta l+x)=-kx$};
\end{tikzpicture}


\figcaption{竖直弹簧振子：以平衡位置为参考，重力项被平衡条件抵消}\label{fig:vertical-spring-shm}
\end{center}

\textbf{证明：} 取平衡位置向下为正方向。
\begin{enumerate}
  \item 平衡时：$k\Delta l = mg$。
  \item 从平衡位置向下拉 $x$：弹力 $F_{\text{弹}} = k(\Delta l + x)$，方向向上。
  \item 合力：
  \[
  F = mg - k(\Delta l + x) = -kx
  \]
  \item 合力与偏离量成正比、方向指向平衡位置，所以是简谐运动。
\end{enumerate}

\textbf{结论：} 竖直弹簧振子也是简谐运动，周期公式 $T = 2\pi\sqrt{m/k}$ 仍然成立——与水平放置时相同。这里的位移 $x$ 必须从新的平衡位置量起，重力只改变平衡位置，不改变回复力系数 $k$。

\textbf{竞赛提示：} 证明简谐运动的关键是找到平衡位置，然后看合力是否与偏离量成正比且指向平衡位置。

\medskip
\textbf{等效弹簧：多个弹簧先化成一个 $k$}

以下推导要求弹簧质量和中间连接件质量可以忽略，弹簧处于弹性限度内，并且只研究共同轴线方向的形变。若物体同时连接多个弹簧，先求等效劲度系数 $k_{\text{等效}}$，再代入
\[
T=2\pi\sqrt{\frac{m}{k_{\text{等效}}}}.
\]

\begin{center}
% TikZ 图：串联、并联弹簧的等效劲度系数
% Day 13：机械振动
\begin{tikzpicture}[scale=0.86,>=Latex,line cap=round,line join=round,
  every node/.style={font=\small}]
  % Parallel springs.
  \begin{scope}[shift={(-3.25,0)}]
    \draw[thick] (-2.05,-0.35) -- (-2.05,1.55);
    \foreach \y in {-0.20,0.20,0.60,1.00,1.40} {
      \draw[gray] (-2.05,\y) -- ++(-0.18,-0.16);
    }
    \draw[decorate,decoration={coil,aspect=0.42,segment length=4pt,amplitude=3pt},thick]
      (-2.05,1.08) -- (0.35,1.08);
    \draw[decorate,decoration={coil,aspect=0.42,segment length=4pt,amplitude=3pt},thick]
      (-2.05,0.18) -- (0.35,0.18);
    \draw[thick,fill=gray!12] (0.35,-0.18) rectangle (1.20,1.45);
    \node at (0.78,0.63) {$m$};
    \node[above] at (-0.85,1.20) {$k_1$};
    \node[below] at (-0.85,0.06) {$k_2$};
    \draw[->,red!80!black,thick] (1.20,0.63) -- (1.85,0.63) node[right] {$x$};
    \node[align=center] at (-0.10,-0.92) {并联：形变量相同\\$k_{\text{并}}=k_1+k_2$};
  \end{scope}

  % Series springs.
  \begin{scope}[shift={(3.25,0)}]
    \draw[thick] (-2.25,0.25) -- (-2.25,1.25);
    \foreach \y in {0.35,0.65,0.95,1.25} {
      \draw[gray] (-2.25,\y) -- ++(-0.18,-0.16);
    }
    \draw[decorate,decoration={coil,aspect=0.42,segment length=4pt,amplitude=3pt},thick]
      (-2.25,0.75) -- (-0.62,0.75);
    \filldraw[black] (-0.62,0.75) circle (1.3pt);
    \draw[decorate,decoration={coil,aspect=0.42,segment length=4pt,amplitude=3pt},thick]
      (-0.62,0.75) -- (1.05,0.75);
    \draw[thick,fill=gray!12] (1.05,0.28) rectangle (1.88,1.22);
    \node at (1.47,0.75) {$m$};
    \node[above] at (-1.50,0.92) {$k_1$};
    \node[above] at (0.25,0.92) {$k_2$};
    \draw[->,red!80!black,thick] (1.88,0.75) -- (2.50,0.75) node[right] {$x$};
    \node[align=center] at (-0.12,-0.92) {串联：拉力相同\\$\dfrac{1}{k_{\text{串}}}=\dfrac{1}{k_1}+\dfrac{1}{k_2}$};
  \end{scope}
\end{tikzpicture}


\figcaption{等效弹簧：并联看力相加，串联看形变量相加}\label{fig:equivalent-springs}
\end{center}

\textbf{并联弹簧：} 两根弹簧一端固定，另一端共同连接同一物体。物体向右偏离 $x$ 时，两根弹簧的形变量都是 $x$，回复力大小分别为 $k_1x$、$k_2x$，总回复力
\[
F=(k_1+k_2)x.
\]
所以
\[
k_{\text{并}}=k_1+k_2.
\]

\textbf{串联弹簧：} 两根弹簧首尾相连。串联时两根弹簧中的拉力大小相同，设为 $F$。它们的形变量分别为
\[
x_1=\frac{F}{k_1},\qquad x_2=\frac{F}{k_2}.
\]
总形变量
\[
x=x_1+x_2=F\left(\frac1{k_1}+\frac1{k_2}\right).
\]
若用一个等效弹簧表示，$x=\dfrac{F}{k_{\text{串}}}$，因此
\[
\frac1{k_{\text{串}}}=\frac1{k_1}+\frac1{k_2}.
\]

\textbf{认知提醒：} 并联弹簧更“硬”，等效 $k$ 变大；串联弹簧更“软”，等效 $k$ 变小。若两根相同弹簧劲度系数都是 $k$，并联得 $2k$，串联得 $\dfrac{k}{2}$。

\medskip
\textbf{参考圆法：用匀速圆周运动理解简谐运动}

简谐运动可以用匀速圆周运动在一条直径上的投影进行数学描述，这称为参考圆法；它不表示振子实际在做圆周运动。若参考圆半径为 $A$，角速度为 $\omega$，从最右端开始计时，则投影满足
\[
x=A\cos\omega t.
\]

\begin{center}
% TikZ 图：参考圆法理解简谐运动
% Day 13：机械振动
\begin{tikzpicture}[scale=0.90,>=Latex]
  \coordinate (O) at (0,0);
  \coordinate (P0) at (0:1.55);
  \coordinate (P1) at (60:1.55);
  \coordinate (H) at ({1.55*cos(60)},0);

  \draw[->,thick] (-2.05,0) -- (2.35,0) node[right] {$x$};
  \draw[dashed,gray] (O) circle (1.55);
  \filldraw[black] (O) circle (1.5pt);
  \node[below left=2pt] at (O) {$O$};
  \draw[->,blue,thick] (20:1.55) arc (20:70:1.55);
  \node[blue] at (44:1.95) {$\omega$};

  \filldraw[red] (P0) circle (2pt);
  \node[above right=2pt] at (P0) {$A$};
  \filldraw[red] (P1) circle (2pt);
  \draw[dashed,red] (P1) -- (H);
  \filldraw[red] (H) circle (2pt);
  \node[below=9pt] at (H) {$\dfrac{A}{2}$};

  \draw[thick] (O) -- (P0);
  \draw[thick,red] (O) -- (P1);
  \draw pic["$\dfrac{\pi}{3}$",draw,->,angle radius=0.55cm,angle eccentricity=1.45] {angle=P0--O--P1};

  \node[align=left,font=\small,anchor=west] at (3.15,0.72)
    {从 $x=A$ 到 $x=A/2$：\\[3pt]
     $\omega t=\dfrac{\pi}{3}$\\[3pt]
     $t=\dfrac{\pi/3}{2\pi/T}=\dfrac{T}{6}$};
\end{tikzpicture}


\figcaption{参考圆法：用匀速圆周运动的投影描述简谐运动的相位变化}\label{fig:reference-circle-shm}
\end{center}

\textbf{例：从最大位移到振幅一半位置要多久？}

若从 $x=A$ 开始计时，到 $x=\dfrac{A}{2}$ 时，
\[
A\cos\omega t=\frac{A}{2},\qquad \cos\omega t=\frac12.
\]
因此
\[
\omega t=\frac{\pi}{3},\qquad
t=\frac{\pi/3}{2\pi/T}=\frac{T}{6}.
\]

\textbf{容易混淆：} 若从平衡位置 $x=0$ 开始向正方向运动，到 $x=\dfrac{A}{2}$，对应参考圆转过角度 $\dfrac{\pi}{6}$，所用时间是 $\dfrac{T}{12}$。所以解题时必须先看清“从哪里开始计时”。

\medskip
\textbf{受迫振动与共振}

真实系统常受到外界周期性驱动，例如秋千被人周期性推动，桥梁受到周期性风力或人群步伐作用。外力频率叫\textbf{驱动频率}，系统自由振动时由自身结构决定的频率叫\textbf{固有频率}。当驱动频率接近固有频率时，系统容易从外界持续获得能量，稳定振幅会显著增大，这就是\textbf{共振}。

\begin{center}
\begin{tikzpicture}[scale=0.95,>=Stealth]
  \draw[->,thick] (0,0) -- (5.4,0) node[right,font=\small] {驱动频率 $f$};
  \draw[->,thick] (0,0) -- (0,3.0) node[above,font=\small] {振幅 $A$};
  \draw[blue,very thick] (0.25,0.35) .. controls (1.15,0.50) and (1.70,1.25) .. (2.45,2.65)
    .. controls (3.00,1.15) and (3.75,0.55) .. (5.05,0.38);
  \draw[dashed] (2.45,0) -- (2.45,2.65);
  \node[font=\small] at (2.45,-0.28) {$f_0$};
  \node[font=\small,red] at (1.92,2.52) {共振};
  \draw[->,red,thick] (2.12,2.42) -- (2.43,2.55);

  \begin{scope}[xshift=0.75cm,yshift=1.95cm]
    \draw[thick] (-0.45,0.30) -- (0.45,0.30);
    \draw[decorate,decoration={coil,aspect=0.45,segment length=3pt,amplitude=3pt},thick] (0,0.30) -- (0,-0.35);
    \draw[thick,fill=gray!15] (-0.18,-0.70) rectangle (0.18,-0.35);
    \draw[<->,orange!85!black,thick] (0.45,-0.72) -- (0.45,-0.18);
  \end{scope}
  \node[font=\scriptsize,align=left,anchor=west] at (3.05,2.42) {驱动频率接近\\固有频率时\\振幅显著增大};
\end{tikzpicture}


\figcaption{共振曲线：驱动频率接近固有频率时振幅显著增大，阻尼越小共振峰越高、越尖}\label{fig:resonance-curve}
\end{center}

\textbf{模型意识：} 判断是否接近共振，先比较驱动频率和系统固有频率，再看阻尼大小。共振不是“外力越大越容易发生”，而是频率接近时外力能够持续有效地向系统输入能量；阻尼越小，共振峰通常越高、越尖。图中把峰值画在 $f_0$ 附近，强调的是中学阶段的近似规律。工程上既会利用共振（乐器、石英振荡器），也要避免危险共振（桥梁、机器、建筑结构）。
\end{competitionbox}

\section*{九、易错点}

\begin{mistakebox}
\begin{enumerate}
  \item \textbf{混淆振幅和位移：} 先找平衡位置。振幅 $A$ 是最大位移的绝对值，是常数；位移 $x$ 是随时间变化的量。

  \item \textbf{单摆公式 $T = 2\pi\sqrt{\dfrac{l}{g}}$ 的适用条件：} 套公式前先检查小角度和单摆条件：摆角较小，摆球直径远小于摆长。

  \item \textbf{弹簧振子中误认为质量不计：} 代入 $T=2\pi\sqrt{m/k}$ 前，先在图上圈出实际振动的物块或小球；它的质量就是式中的 $m$，不能省去。题设“轻弹簧”仅表示弹簧自身质量可忽略，并不表示振子质量为零。

  \item \textbf{回复力不是新力：} 先画真实力，再看合力或分力是否指向平衡位置；回复力是效果力，不是一种新力。

  \item \textbf{单摆周期与质量无关：} 先检查是不是理想单摆、小角度摆动；条件满足时，公式中没有 $m$，周期与质量无关。若摆角很大、摆线有质量或空气阻力明显，就不能只凭这个结论判断。

  \item \textbf{把串联、并联弹簧公式记反：} 并联弹簧形变量相同、弹力相加；串联弹簧拉力相同、形变量相加。
\end{enumerate}
\end{mistakebox}

\section*{十、今日小结}

\begin{enumerate}
  \item 简谐运动回复力 $F=-Kx$；对弹簧振子，回复系数 $K$ 就是劲度系数 $k$。
  \item 理想轻弹簧振子周期 $T = 2\pi\sqrt{\dfrac{m}{k}}$，在线性范围内与振幅无关。
  \item 单摆周期 $T = 2\pi\sqrt{\dfrac{l}{g}}$，理想小角度近似下与质量无关，并近似与振幅无关。
  \item 简谐运动 $x$-$t$ 图像为正弦或余弦曲线，$x=A\cos(\omega t+\varphi_0)$，且 $v_{\max}=\omega A$、$a=-\omega^2x$。
  \item 理想无阻尼简谐运动中机械能守恒，弹簧振子 $E = \frac12 kA^2$。
  \item 多根弹簧先求等效劲度系数：并联 $k$ 相加，串联倒数相加。
  \item 共振发生在驱动频率接近固有频率时，关键是频率匹配。
\end{enumerate}

\textbf{一分钟复盘：} 任选弹簧振子或单摆，画一张小图并标出平衡位置、位移 $x$ 的量法和回复力方向；再在图旁写一行“周期公式 + 适用条件”。若平衡位置和位移起点说不清，先不要代公式。
